\begin{afterword}
\summaryhead

The post starts from Feynman's question: what poets can speak of Jupiter as a man, but must be silent if he is a spinning sphere of methane and ammonia? The author gives a real answer. A Jupiter like us can love and weep, and the strongest human emotions, found in every culture on Donald Brown's list of human universals, are built into us. A sphere of gas is harder to write about. So Keats felt a loss when the rainbow became raindrops, and would feel it even knowing the mathematics, since the Biblical story of the flood is more dramatic than refraction.

The author disagrees. The Great Stories do not need Jupiter, because humans still love, hate and kill every day. Nor are they timeless: the species that tells them is not. The author hopes some of them will change, and asks poets to tell the story of strange new things truly, as James Thomson did for Newton's rainbow. Poets who can tell only the old campfire stories are ``savanna poets.''

\responsehead

The post builds the strongest case it can for the other side and says that it did so. The Romantic poet's best reply, that equations are not about strong emotions, is the author's own construction, offered as ``the strongest rejoinder I can think of.'' The answer, that humans remain to carry the stories, meets that reply directly. Its best example is Thomson's poem of 1727, which describes the colours ``adown the watery bow'' and asks whether any poet ever imagined anything so fair, about ninety years before \textsc{Lamia}.

Two steps are asserted rather than argued. The first is the claim that the strongest emotions are ``deeply engraved, blood and bone, brain and DNA.'' Brown's list shows that such things are found in every culture, not why, and the same list includes weapons and units of time. The post meets possible doubt by saying that no one who knows evolutionary psychology would deny it. The argument about poets would stand on universality alone; the claim about DNA goes further, and is given without evidence.

The second is the case against timelessness. The poets' claim, as the post reports it, is that Sophocles and Shakespeare could have swapped centuries. The post answers that far enough back in hominid evolution no one would understand \textsc{Hamlet}. That is true of ``timeless'' taken literally, but over the span the poets meant, the post agrees with them. The real disagreement is about the future, and there the post states a hope, and a forecast offered as a sincere doubt that we have ``another thousand years to go in our current form,'' with no argument in this post.

What remains is an aesthetic program, not a finding: stories should diversify, and poets should write about what is new and true. The post argues for it by example and preference, which is what such a program allows.

In short: a fair-minded answer to Feynman's question that states the other side at its strongest, with an unargued claim about emotions and DNA and a rebuttal of ``timeless'' that concedes the poets' point for the past.
\end{afterword}
